% 第11章第8节：伸缩分离与正交逼近（开方方程、正定平方根、极分解、最佳正交逼近、对称矩阵的谱函数）。

\section{伸缩分离与正交逼近}
\label{sec:ls-polar-decomposition}

复数$z\neq0$可以写成$z=e^{\mathrm i\theta}r$：先乘正数$r$改变长度，再乘单位复数$e^{\mathrm i\theta}$转动方向。本节对实方阵提出同样的问题：给定$\bm A\in\mathbb R^{d\times d}$，怎样求出$\bm A=\bm W\bm H$，使$\bm H$沿一组互相垂直的方向作非负伸缩，$\bm W$保持长度？

两个因子都未知，不能直接分别求解。我们先处理可逆的$\bm A$：利用$\bm W^\top\bm W=\bm I$消去未知的$\bm W$，得到只含$\bm H$的方程$\bm H^2=\bm A^\top\bm A$；再由已知矩阵$\bm A^\top\bm A$的谱分解开方构造$\bm H$，并通过$\bm W=\bm A\bm H^{-1}$构造和验证$\bm W$。两个因子求出以后，再处理奇异情形，并说明所得的$\bm W$也是离$\bm A$最近的正交矩阵。最后，把开方与求逆中共同使用的“沿谱改写特征值”归纳为对称矩阵的函数。

\subsection{消去正交因子，得到开方方程}
\label{sec:ls-polar-square-root-equation}

先看一个具体矩阵。取
\[
 \bm A=\begin{bmatrix}0&-1\\3&0\end{bmatrix},
\]
它把$\bm e_1$送到$3\bm e_2$，把$\bm e_2$送到$-\bm e_1$。这个作用既改变长度，又改变方向：$\bm e_1$被拉长为原来的$3$倍，同时两个坐标轴都被转了$90$度。我们希望把它拆成两步：先由对称正定的$\bm H$沿一组正交方向伸缩，再由正交矩阵$\bm W$保长地改变方向，即求$\bm A=\bm W\bm H$。

本小节和下一小节先假定$\bm A$可逆。这时$\bm A$不把任何非零向量压成零，伸缩可以撤销，所求的$\bm H$应当正定。正交矩阵既可以是旋转，也可以包含反射，这里不加限制；$\bm A$奇异的情形在第~\ref{sec:ls-polar-singular-best-orthogonal}节处理。

两个因子同时未知，先设法消去一个。若$\bm A=\bm W\bm H$，其中$\bm W^\top\bm W=\bm I$、$\bm H^\top=\bm H$，则把等式两边分别转置后与原式相乘，得到
\begin{equation}
 \bm A^\top\bm A
 =\bm H^\top\bm W^\top\bm W\bm H
 =\bm H^2 .
 \label{eq:ls-polar-square-root-equation}
\end{equation}
正交因子在相乘时消失了。记$\bm S=\bm A^\top\bm A$，它完全由已知的$\bm A$确定，于是原问题的前一半化为一个开方问题：找一个对称正定矩阵$\bm H$，使$\bm H^2=\bm S$。需要注意，式~\eqref{eq:ls-polar-square-root-equation}只是分解若存在就必须满足的条件；这样的$\bm H$是否存在、由它能否再得到正交的$\bm W$，都还没有证明。

开方之前先确认$\bm S$的结构。它对称，因为$(\bm A^\top\bm A)^\top=\bm A^\top\bm A$；对任意非零$\bm x$，
\[
 \bm x^\top\bm S\bm x=(\bm A\bm x)^\top(\bm A\bm x)=\lVert\bm A\bm x\rVert_2^2>0,
\]
最后一步用到$\bm A$可逆，从而$\bm A\bm x\neq\bm0$。所以$\bm S$对称正定，可以使用第~\ref{sec:ls-spectral-theory}节的正交谱分解，且全部特征值为正。

在贯穿例中，
\[
 \bm S=\bm A^\top\bm A
 =\begin{bmatrix}0&3\\-1&0\end{bmatrix}
 \begin{bmatrix}0&-1\\3&0\end{bmatrix}
 =\begin{bmatrix}9&0\\0&1\end{bmatrix}.
\]
这里$\bm S$恰好是对角矩阵，最直接的做法是沿两个坐标轴分别取正平方根，得到候选$\bm H=\operatorname{diag}(3,1)$。一般的$\bm S$不是对角矩阵，下一小节用谱分解把它化到这种情形。

\subsection{由谱分解构造伸缩因子}
\label{sec:ls-matrix-square-root}

对上一小节的$\bm S=\bm A^\top\bm A$作正交谱分解$\bm S=\bm Q\bm\Lambda\bm Q^\top$，其中$\bm Q=[\bm q_1,\ldots,\bm q_d]$正交，$\bm\Lambda=\operatorname{diag}(\lambda_1,\ldots,\lambda_d)$，全部$\lambda_i>0$。在特征方向$\bm q_i$上，$\bm S$只作标量伸缩$\lambda_i$倍。既然希望$\bm H$作用两次等于$\bm S$，自然尝试让$\bm H$沿同一方向伸缩$\sqrt{\lambda_i}$倍，即构造
\begin{equation}
 \bm H=\bm Q\operatorname{diag}\bigl(\sqrt{\lambda_1},\ldots,\sqrt{\lambda_d}\bigr)\bm Q^\top .
 \label{eq:ls-spectral-square-root}
\end{equation}
这是我们主动选取的一个候选，并没有预先假定未知的$\bm H$必须与$\bm S$共用特征基；它是否满足要求，要靠验证。

验证只需矩阵乘法。记$\bm\Lambda^{1/2}=\operatorname{diag}(\sqrt{\lambda_1},\ldots,\sqrt{\lambda_d})$，由$\bm Q^\top\bm Q=\bm I$，
\[
 \bm H^2
 =\bm Q\bm\Lambda^{1/2}(\bm Q^\top\bm Q)\bm\Lambda^{1/2}\bm Q^\top
 =\bm Q\bm\Lambda\bm Q^\top
 =\bm S .
\]
又$\bm H^\top=\bm Q(\bm\Lambda^{1/2})^\top\bm Q^\top=\bm H$，而式~\eqref{eq:ls-spectral-square-root}本身就是$\bm H$的正交谱分解，特征值$\sqrt{\lambda_i}$全为正，所以$\bm H$对称正定。所需的伸缩因子已经求出，记作$\bm S^{1/2}$，称为$\bm S$的正定平方根，也称\term{矩阵平方根}（matrix square root）。

构造本身只用到$\lambda_i\geq0$，所以对半正定的$\bm S$同样可行，此时得到半正定的$\bm H$。还剩一个问题：$\bm Q$不唯一，特征值有重根时，重根空间内可以换用别的标准正交基；方程$\bm X^2=\bm S$也可能有许多解，例如$\bm I$与$-\bm I$的平方都是$\bm I$。下面的定理说明，只要要求解对称半正定，答案就唯一，因而与$\bm Q$的选取无关。

\begin{theorem}[半正定平方根的存在与唯一]
\label{thm:ls-psd-square-root}
设$\bm S\in\mathbb R^{d\times d}$对称半正定，有正交谱分解$\bm S=\bm Q\bm\Lambda\bm Q^\top$。则存在唯一的对称半正定矩阵$\bm X$满足$\bm X^2=\bm S$，它就是式~\eqref{eq:ls-spectral-square-root}给出的$\bm S^{1/2}$。$\bm S$正定时，$\bm S^{1/2}$也正定。
\end{theorem}
\begin{proof}
存在性与正定性已由上面的验证给出：$\bm S$半正定时全部$\lambda_i\geq0$，$\sqrt{\lambda_i}$有定义，式~\eqref{eq:ls-spectral-square-root}是对称半正定矩阵且平方为$\bm S$；$\bm S$正定时全部$\sqrt{\lambda_i}>0$。

下证唯一性。设$\bm X$对称半正定且$\bm X^2=\bm S$。对$\bm X$本身作正交谱分解$\bm X=\sum_{j=1}^d\mu_j\bm p_j\bm p_j^\top$，其中$\bm p_1,\ldots,\bm p_d$是标准正交基，$\mu_j\geq0$。由$\bm p_j^\top\bm p_k=\delta_{j,k}$，
\[
 \bm S=\bm X^2=\sum_{j=1}^d\mu_j^2\bm p_j\bm p_j^\top ,
\]
所以$\bm S\bm p_j=\mu_j^2\bm p_j$，每个$\bm p_j$都是$\bm S$的特征向量。

取$\bm S$的任一特征值$\lambda$，考察特征空间$\ker(\bm S-\lambda\bm I)$。把任意$\bm x$展开为$\bm x=\sum_jc_j\bm p_j$，则
\[
 (\bm S-\lambda\bm I)\bm x=\sum_{j=1}^dc_j(\mu_j^2-\lambda)\bm p_j ,
\]
由$\bm p_j$线性无关，它为零当且仅当凡$\mu_j^2\neq\lambda$的$j$都有$c_j=0$。因此$\ker(\bm S-\lambda\bm I)$恰由满足$\mu_j^2=\lambda$的那些$\bm p_j$张成。对这些$j$，$\mu_j\geq0$迫使$\mu_j=\sqrt\lambda$，于是对特征空间中的每个$\bm x$，
\[
 \bm X\bm x=\sum_{\mu_j^2=\lambda}c_j\mu_j\bm p_j=\sqrt\lambda\,\bm x .
\]
另一方面，由式~\eqref{eq:ls-spectral-square-root}，$\bm S^{1/2}\bm q_i=\sqrt{\lambda_i}\bm q_i$，而$\ker(\bm S-\lambda\bm I)$由满足$\lambda_i=\lambda$的$\bm q_i$张成，所以$\bm S^{1/2}$在这个特征空间上同样是乘以$\sqrt\lambda$。

这样，$\bm X$与$\bm S^{1/2}$在$\bm S$的每个特征空间上作用相同。由谱定理，这些特征空间合起来张成$\mathbb R^d$（各$\bm q_i$都落在其中某一个里），线性映射在一组基上相同即处处相同，故$\bm X=\bm S^{1/2}$。
\end{proof}

唯一性的论证直接说明了两件事。第一，在重根空间内换基，相应的$\bm q_i$变了，但特征空间和其上的作用“乘以$\sqrt\lambda$”不变，所以$\bm S^{1/2}$不依赖$\bm Q$的选取。第二，唯一性只针对半正定解；去掉这一要求，$\bm X^2=\bm S$一般有许多解，$-\bm S^{1/2}$就是其中之一。

在贯穿例中，$\bm S=\operatorname{diag}(9,1)$，可取$\bm Q=\bm I$，得到$\bm H=\operatorname{diag}(3,1)$，直接核对$\bm H^2=\operatorname{diag}(9,1)=\bm S$。上一小节凭直觉写出的候选正是唯一的正定平方根。到这里，解题进度是：已经得到伸缩因子$\bm H$，还需要求出$\bm W$并验证它正交。这里所做的“把每个特征值换成它的平方根”，将在第~\ref{sec:ls-matrix-functions}节成为矩阵函数的一个具体例子；但构造和验证本身只用到谱分解、标量开方与矩阵乘法。

所得的平方根带来一个后文直接使用的长度恒等式。由$\bm S^{1/2}$对称及$(\bm S^{1/2})^2=\bm S$，对半正定$\bm S$与任意$\bm x$，
\begin{equation}
 \lVert\bm S^{1/2}\bm x\rVert_2^2
 =\bm x^\top\bm S^{1/2}\bm S^{1/2}\bm x
 =\bm x^\top\bm S\bm x .
 \label{eq:ls-square-root-length}
\end{equation}
它把第~\ref{sec:ls-weighted-inner-products}节以$\bm S$加权的长度化为普通欧氏长度。对$\bm S=\bm A^\top\bm A$，它给出$\lVert\bm H\bm x\rVert_2=\lVert\bm A\bm x\rVert_2$，说明$\bm H$承担了$\bm A$的全部长度变化；但这并不表示$\bm H\bm x$与$\bm A\bm x$方向相同，方向的差别正要由$\bm W$补上。第~\ref{sec:ls-stationary-iterations}节分析Gauss--Seidel迭代的收缩时，就用式~\eqref{eq:ls-square-root-length}把能量范数写成欧氏长度。

若只需表示同一个长度，平方根并不是唯一的选择。由定理~\ref{thm:ls-gram-factor}，正定$\bm S$还可写成$\bm S=\bm L^\top\bm L$，$\bm L$可逆，此时同样有$\lVert\bm L\bm x\rVert_2^2=\bm x^\top\bm S\bm x=\lVert\bm S^{1/2}\bm x\rVert_2^2$。不过这只说明两个向量$\bm L\bm x$与$\bm S^{1/2}\bm x$长度相同，并不表示两个矩阵对向量的作用相同：$\bm L$一般不对称，$\bm S^{1/2}$则由定理~\ref{thm:ls-psd-square-root}唯一确定。下一小节会看到，二者恰好相差一个正交因子。第~\ref{chap:linear-equations-and-optimization}章的Cholesky分解把$\bm L$取成三角矩阵，那时作用$\bm L^{-1}$只需解三角方程，不必求出完整的谱分解。

\subsection{由伸缩因子构造正交因子}
\label{sec:ls-polar-orthogonal-factor}

回到$\bm A=\bm W\bm H$。刚求出的$\bm H$正定，因而可逆，在等式两边右乘$\bm H^{-1}$，得到正交因子唯一可能的候选$\bm W=\bm A\bm H^{-1}$。右乘而不是左乘，是由原方程中因子的顺序决定的：$\bm H$先作用，要撤销它的伸缩，就要在$\bm A$的右侧作用$\bm H^{-1}$。

$\bm H^{-1}$同样可以沿特征方向写出。$\bm H$沿$\bm q_i$伸缩$\sqrt{\lambda_i}$倍，撤销这一伸缩就是沿同一方向乘以$1/\sqrt{\lambda_i}$。令
\[
 \bm H^{-1}=\bm Q\operatorname{diag}\bigl(1/\sqrt{\lambda_1},\ldots,1/\sqrt{\lambda_d}\bigr)\bm Q^\top ,
\]
由$\bm Q^\top\bm Q=\bm I$，它与$\bm H$相乘得$\bm Q\bm I\bm Q^\top=\bm I$，确实是逆矩阵。它也对称正定，记作$\bm S^{-1/2}$，称为$\bm S$的逆平方根。

定义了$\bm W=\bm A\bm H^{-1}$之后，必须检查它满足原来要求的性质。由$(\bm H^{-1})^\top=(\bm H^\top)^{-1}=\bm H^{-1}$及$\bm A^\top\bm A=\bm S=\bm H^2$，
\[
 \bm W^\top\bm W
 =\bm H^{-1}\bm A^\top\bm A\bm H^{-1}
 =\bm H^{-1}\bm H^2\bm H^{-1}
 =\bm I ,
\]
所以$\bm W$正交；又$\bm W\bm H=\bm A\bm H^{-1}\bm H=\bm A$。正交性和分解等式都成立，存在性才算证完。这样的分解有一个专门的名称。

\begin{definition}[极分解]
\label{def:matrix-analysis-polar-decomposition}
设$\bm A\in\mathbb R^{d\times d}$。若$\bm W$是正交矩阵，$\bm H$是对称半正定矩阵，且
\begin{equation}
\begin{gathered}
 \bm A=\bm W\bm H\\
 \bm W^\top\bm W=\bm I,\quad \bm H=(\bm A^\top\bm A)^{1/2},
\end{gathered}
\label{eq:matrix-analysis-polar-decomposition}
\end{equation}
则称$\bm A=\bm W\bm H$为$\bm A$的\term{极分解}（polar decomposition），$\bm W$称为正交因子，$\bm H$称为伸缩因子。
\end{definition}

式中$\bm H=(\bm A^\top\bm A)^{1/2}$并不是额外的要求：由式~\eqref{eq:ls-polar-square-root-equation}，任何满足$\bm A=\bm W\bm H$、$\bm W$正交、$\bm H$对称半正定的分解都有$\bm H^2=\bm A^\top\bm A$，再由定理~\ref{thm:ls-psd-square-root}，$\bm H$只能是$(\bm A^\top\bm A)^{1/2}$。把上面的推导与验证合在一起，就得到可逆情形的结论。

\begin{theorem}[可逆矩阵的极分解]
\label{thm:ls-polar-invertible}
设$\bm A\in\mathbb R^{d\times d}$可逆，则$\bm A$有极分解，且两个因子都唯一：
\[
 \bm H=(\bm A^\top\bm A)^{1/2},\qquad
 \bm W=\bm A(\bm A^\top\bm A)^{-1/2},
\]
其中$\bm H$对称正定。
\end{theorem}
\begin{proof}
存在性：由第~\ref{sec:ls-polar-square-root-equation}节，$\bm S=\bm A^\top\bm A$对称正定；由定理~\ref{thm:ls-psd-square-root}，$\bm H=\bm S^{1/2}$对称正定，从而可逆；上面已验证$\bm W=\bm A\bm H^{-1}$正交且$\bm W\bm H=\bm A$。

唯一性：设$\bm A=\bm W'\bm H'$是任一极分解。由式~\eqref{eq:ls-polar-square-root-equation}，$\bm H'^2=\bm S$，又$\bm H'$对称半正定，由定理~\ref{thm:ls-psd-square-root}得$\bm H'=\bm S^{1/2}=\bm H$。$\bm H$可逆，在$\bm A=\bm W'\bm H$两边右乘$\bm H^{-1}$得$\bm W'=\bm A\bm H^{-1}=\bm W$。
\end{proof}

由$\det\bm A=\det\bm W\det\bm H$及$\det\bm H>0$，$\det\bm W$与$\det\bm A$同号：$\det\bm A>0$时$\bm W$是旋转，$\det\bm A<0$时$\bm W$包含一次反射。这正是第~\ref{sec:ls-polar-square-root-equation}节不限制正交作用类型的原因。上一小节的Gram因子也可以由此理解：对可逆的$\bm L$，$\bm L^\top\bm L=\bm S$，所以$\bm L$的极分解是$\bm L=(\bm L\bm S^{-1/2})\bm S^{1/2}$，$\bm L$与$\bm S^{1/2}$恰好相差正交因子$\bm L\bm S^{-1/2}$，这解释了二者为何给出相同的长度。

现在把贯穿例算完。$\bm H=\operatorname{diag}(3,1)$，所以$\bm H^{-1}=\operatorname{diag}(1/3,1)$，右乘得
\[
 \bm W=\bm A\bm H^{-1}
 =\begin{bmatrix}0&-1\\3&0\end{bmatrix}
 \begin{bmatrix}1/3&0\\0&1\end{bmatrix}
 =\begin{bmatrix}0&-1\\1&0\end{bmatrix}.
\]
逐项核对：$\bm W^\top\bm W=\bm I$；$\bm W\bm H=\left[\begin{smallmatrix}0&-1\\3&0\end{smallmatrix}\right]=\bm A$。$\bm W$把$\bm e_1$送到$\bm e_2$、把$\bm e_2$送到$-\bm e_1$，是逆时针旋转$90$度。于是原变换可以读成：先沿$\bm e_1$伸长$3$倍、沿$\bm e_2$保持不变，再整体旋转$90$度。在这个解题过程中，$\bm S$是由$\bm A$直接算出的中间量，$\bm H$是对$\bm S$开方得到的因子，$\bm W$则是撤销伸缩后剩下的因子。

\subsection{奇异情形与最佳正交逼近}
\label{sec:ls-polar-singular-best-orthogonal}

\paragraph{零伸缩方向上的正交补全}

当$\bm A$奇异时，$\bm S=\bm A^\top\bm A$仍对称半正定，但有零特征值。定理~\ref{thm:ls-psd-square-root}对半正定矩阵成立，所以沿谱取非负平方根仍能构造唯一的半正定$\bm H=\bm S^{1/2}$。上一小节的做法却在下一步失效：零特征值不能取倒数，$\bm H^{-1}$不存在，也就无法由$\bm W=\bm A\bm H^{-1}$得到正交因子。

这时可以直接利用第~\ref{sec:ls-svd}节已有的奇异值分解。由定理~\ref{thm:ls-svd-existence}，方阵$\bm A$可以写成$\bm A=\bm U\bm\Sigma\bm V^\top$，其中$\bm U,\bm V\in\mathbb R^{d\times d}$正交，$\bm\Sigma=\operatorname{diag}(\sigma_1,\ldots,\sigma_d)$，$\sigma_1\geq\cdots\geq\sigma_d\geq0$。在$\bm U$与$\bm\Sigma$之间插入$\bm V^\top\bm V=\bm I$，就得到候选
\begin{equation}
\begin{gathered}
 \bm A=(\bm U\bm V^\top)(\bm V\bm\Sigma\bm V^\top)\\
 \bm W=\bm U\bm V^\top,\quad \bm H=\bm V\bm\Sigma\bm V^\top .
\end{gathered}
\label{eq:ls-polar-from-svd}
\end{equation}

\begin{theorem}[一般方阵的极分解]
\label{thm:ls-polar-general}
每个$\bm A\in\mathbb R^{d\times d}$都有极分解，式~\eqref{eq:ls-polar-from-svd}给出其中一个。伸缩因子$\bm H=(\bm A^\top\bm A)^{1/2}$总是唯一的；正交因子$\bm W$唯一，当且仅当$\bm A$可逆。
\end{theorem}
\begin{proof}
先验证式~\eqref{eq:ls-polar-from-svd}。$\bm H=\bm V\bm\Sigma\bm V^\top$是对称矩阵的正交谱分解，特征值$\sigma_i\geq0$，故$\bm H$对称半正定；由$\bm U^\top\bm U=\bm I$与$\bm\Sigma^\top=\bm\Sigma$，
\[
 \bm H^2=\bm V\bm\Sigma^2\bm V^\top
 =\bm V\bm\Sigma^\top\bm U^\top\bm U\bm\Sigma\bm V^\top
 =\bm A^\top\bm A .
\]
又$(\bm U\bm V^\top)^\top\bm U\bm V^\top=\bm V\bm U^\top\bm U\bm V^\top=\bm I$，所以$\bm W$正交；而$\bm W\bm H=\bm U\bm V^\top\bm V\bm\Sigma\bm V^\top=\bm U\bm\Sigma\bm V^\top=\bm A$。这就证明了存在性。

$\bm H$的唯一性与可逆情形相同：任何极分解都满足式~\eqref{eq:ls-polar-square-root-equation}，由定理~\ref{thm:ls-psd-square-root}，$\bm H=(\bm A^\top\bm A)^{1/2}$。$\bm A$可逆时，$\bm W$的唯一性已由定理~\ref{thm:ls-polar-invertible}给出。

设$\bm A$奇异，则$\operatorname{rank}\bm A<d$，由定理~\ref{thm:ls-svd-existence}知$\sigma_d=0$。令$\bm U'$为把$\bm U$的最后一列$\bm u_d$换成$-\bm u_d$所得的矩阵，它仍正交。$\bm U\bm\Sigma$的第$i$列是$\sigma_i\bm u_i$，最后一列因$\sigma_d=0$而为零，所以$\bm U'\bm\Sigma=\bm U\bm\Sigma$，从而$\bm A=\bm U'\bm\Sigma\bm V^\top$。按同样的验证，$\bm W'=\bm U'\bm V^\top$与同一个$\bm H$也构成极分解，而
\[
 \bm W-\bm W'=(\bm U-\bm U')\bm V^\top=2\bm u_d\bm v_d^\top\neq\bm0 ,
\]
所以正交因子不唯一。
\end{proof}

不唯一的原因可以从作用上看清。设$r=\operatorname{rank}\bm A$。对$i\leq r$，$\bm H\bm v_i=\sigma_i\bm v_i$且$\sigma_i>0$，由$\bm A=\bm W\bm H$得$\bm W\bm v_i=\bm A\bm v_i/\sigma_i=\bm u_i$，正交因子在这些方向上由$\bm A$完全确定。对$i>r$，$\bm H$把$\bm v_i$压成零，$\bm A$在这些方向上不提供任何信息，$\bm W$只需把它们标准正交地送到$\bm u_1,\ldots,\bm u_r$张成空间的正交补中；这一步是正交补全，可以有多种选择。这里始终没有定义不存在的$\bm H^{-1}$，也不能用逐元素取倒数代替矩阵求逆。

可逆时，式~\eqref{eq:ls-polar-from-svd}与定理~\ref{thm:ls-polar-invertible}给出的是同一对因子，因为两者都是极分解，而可逆情形的极分解唯一。谱开方的路线解释了因子为什么是这个样子；实际计算时，若已经有SVD，可以直接读出两个因子。

\paragraph{正交因子的最优性}

构造完成后，可以进一步比较全部正交矩阵：极分解的正交因子是不是Frobenius范数下离$\bm A$最近的正交矩阵？为避免与谱分解的$\bm Q$混淆，下面用$\bm Z$表示任意候选正交矩阵。由$\lVert\bm Z\rVert_{\mathrm F}^2=\operatorname{tr}(\bm Z^\top\bm Z)=\operatorname{tr}\bm I=d$，
\[
 \lVert\bm A-\bm Z\rVert_{\mathrm F}^2
 =\lVert\bm A\rVert_{\mathrm F}^2+d-2\operatorname{tr}(\bm Z^\top\bm A),
\]
前两项与$\bm Z$无关，所以只需在正交矩阵中最大化$\operatorname{tr}(\bm Z^\top\bm A)$。这个目标关于$\bm Z$是线性的；第~\ref{sec:ls-rank-k-trace}节选择子空间时的$\operatorname{tr}(\bm V^\top\bm G\bm V)$关于$\bm V$是二次的，那里的迹极值界不能直接搬用，需要重新给出上界。

\begin{theorem}[最佳正交逼近]
\label{thm:ls-best-orthogonal-approximation}
设$\bm A\in\mathbb R^{d\times d}$有SVD$\bm A=\bm U\bm\Sigma\bm V^\top$。对任意正交矩阵$\bm Z$，
\[
 \operatorname{tr}(\bm Z^\top\bm A)\leq\sum_{i=1}^d\sigma_i ,
\]
取$\bm Z=\bm U\bm V^\top$时等号成立。因此
\[
 \min_{\bm Z^\top\bm Z=\bm I}\lVert\bm A-\bm Z\rVert_{\mathrm F}^2
 =\sum_{i=1}^d(\sigma_i-1)^2 ,
\]
最小值由极分解的正交因子达到。$\bm A$可逆时最优解唯一，就是$\bm W=\bm A(\bm A^\top\bm A)^{-1/2}$。
\end{theorem}
\begin{proof}
令$\bm C=\bm U^\top\bm Z\bm V$，它是正交矩阵的乘积，仍然正交，且$\bm Z=\bm U\bm C\bm V^\top$。由迹的循环性，
\[
 \operatorname{tr}(\bm Z^\top\bm A)
 =\operatorname{tr}(\bm V\bm C^\top\bm U^\top\bm U\bm\Sigma\bm V^\top)
 =\operatorname{tr}(\bm C^\top\bm\Sigma)
 =\sum_{i=1}^d\sigma_iC_{i,i}.
\]
$\bm C$的每一列都是单位向量，所以$|C_{i,i}|\leq1$；又$\sigma_i\geq0$，故$\sigma_iC_{i,i}\leq\sigma_i$，求和即得上界。$\bm Z=\bm U\bm V^\top$时$\bm C=\bm I$，等号成立。

由式~\eqref{eq:matrix-analysis-frobenius-norm}，$\lVert\bm A\rVert_{\mathrm F}^2=\sum_i\sigma_i^2$，代入展开式得最小值
\[
 \sum_{i=1}^d\sigma_i^2+d-2\sum_{i=1}^d\sigma_i
 =\sum_{i=1}^d(\sigma_i-1)^2 .
\]
由定理~\ref{thm:ls-polar-general}，$\bm U\bm V^\top$是极分解的正交因子。

设$\bm A$可逆，则全部$\sigma_i>0$。等号要求每一项$\sigma_iC_{i,i}=\sigma_i$，即$C_{i,i}=1$；单位列向量$\bm C_{:,i}$的第$i$个分量为$1$，其余分量只能为零，于是$\bm C_{:,i}=\bm e_i$，$\bm C=\bm I$，$\bm Z=\bm U\bm V^\top$。由定理~\ref{thm:ls-polar-invertible}，它就是$\bm W$。
\end{proof}

去掉伸缩得到的正交因子，同时解决了“离$\bm A$最近的正交矩阵”这个问题。含噪的旋转估计常用这一步，把一个近似正交的矩阵修正为正交矩阵。在贯穿例中，$\bm S=\operatorname{diag}(9,1)$，奇异值为$3$与$1$，最小平方误差为$(3-1)^2+(1-1)^2=4$；直接计算也得$\bm A-\bm W=\left[\begin{smallmatrix}0&0\\2&0\end{smallmatrix}\right]$，其Frobenius范数的平方为$4$。

$\bm A$奇异时，最优解一般不唯一，而且最优解恰好是全部极分解的正交因子。一方面，若$\bm A=\bm W'\bm H$是任一极分解，则$\operatorname{tr}(\bm W'^\top\bm A)=\operatorname{tr}(\bm W'^\top\bm W'\bm H)=\operatorname{tr}\bm H=\sum_i\sigma_i$，达到上界。另一方面，设$r=\operatorname{rank}\bm A$，$\sigma_i=0$（$i>r$）的项对迹没有贡献，等号只要求$C_{i,i}=1$（$i\leq r$）。按上面的论证，$\bm C$的前$r$列是$\bm e_1,\ldots,\bm e_r$；由正交性，其余各列与它们正交，于是$\bm C=\operatorname{diag}(\bm I_r,\bm C')$，其中$\bm C'$是某个$d-r$阶正交矩阵。$\bm U\bm C$仍正交，且只改变了$\bm U$的后$d-r$列，这些列在$\bm U\bm\Sigma$中都乘以零，所以$\bm A=(\bm U\bm C)\bm\Sigma\bm V^\top$也是SVD，由定理~\ref{thm:ls-polar-general}，$\bm Z=(\bm U\bm C)\bm V^\top$是极分解的正交因子。不同的$\bm C'$正对应前面所说零伸缩方向上不同的正交补全。

最后说明这一结论的边界。最优解是在全部正交矩阵中取得的；$\bm A$可逆时，唯一的最优解$\bm W$满足$\det\bm W$与$\det\bm A$同号（第~\ref{sec:ls-polar-orthogonal-factor}节）。若问题额外要求$\bm Z$是旋转，即$\det\bm Z=1$，而$\det\bm A<0$，则$\bm W$不满足约束，带约束的最优解需要另行处理，本节不展开。贯穿例中$\det\bm A=3>0$，$\bm W$本身就是旋转。

\subsection{从具体构造到矩阵函数}
\label{sec:ls-matrix-functions}
\label{sec:matrix-analysis-functions-structured-operations}

\paragraph{沿谱运算的共同方法}

回看前面两次构造。求$\bm H$时，我们保持$\bm S$的特征方向，把每个$\lambda_i$换成$\sqrt{\lambda_i}$；求$\bm H^{-1}$时，同样保持特征方向，把$\lambda_i$换成$1/\sqrt{\lambda_i}$。两次都是“方向不动，只对特征值作用一个标量函数”。把这个实际用过的操作归纳出来，就得到对称矩阵的函数。

设实对称矩阵$\bm S$有正交谱分解$\bm S=\bm Q\bm\Lambda\bm Q^\top$。按第~\ref{sec:ls-spectral-projections}节，把同一特征值的项合并，可以写成$\bm S=\sum_\lambda\lambda\bm P_\lambda$，求和遍历$\bm S$的不同特征值，$\bm P_\lambda$是到特征空间$\ker(\bm S-\lambda\bm I)$的正交投影，满足$\bm P_\lambda\bm P_\mu=\bm0$（$\lambda\neq\mu$）、$\bm P_\lambda^2=\bm P_\lambda$及$\sum_\lambda\bm P_\lambda=\bm I$。

\begin{definition}[对称矩阵的谱函数]
\label{def:ls-symmetric-matrix-function}
设$\bm S\in\mathbb R^{d\times d}$实对称，$f$是定义在$\bm S$全部特征值上的实函数。令
\begin{equation}
\begin{aligned}
 f(\bm S)
 &=\bm Q\operatorname{diag}\bigl(f(\lambda_1),\ldots,f(\lambda_d)\bigr)\bm Q^\top\\
 &=\sum_\lambda f(\lambda)\bm P_\lambda ,
\end{aligned}
\label{eq:ls-symmetric-spectral-function}
\end{equation}
称为由$f$沿谱定义的\term{矩阵函数}（matrix function）。
\end{definition}

两种写法相等：按列展开，第一式是$\sum_if(\lambda_i)\bm q_i\bm q_i^\top$，把$\lambda_i=\lambda$的项合并，$\sum_{\lambda_i=\lambda}\bm q_i\bm q_i^\top=\bm P_\lambda$，就得到第二式。第二式只涉及特征值和谱投影，而谱投影由$\bm S$唯一确定，所以$f(\bm S)$与重根空间内基的选取无关，定义是良好的。$f(z)=\sqrt z$时，它回到第~\ref{sec:ls-matrix-square-root}节的$\bm S^{1/2}$；$f(z)=1/\sqrt z$时，回到第~\ref{sec:ls-polar-orthogonal-factor}节的$\bm S^{-1/2}$。前面那些构造在这里得到总结，而不是反过来，必须先建立一般的矩阵函数才能解出$\bm H$与$\bm W$。

推广之后要检查它与已有的矩阵运算相容，否则同一个记号会有两种含义。

\begin{lemma}[谱函数与矩阵运算相容]
\label{lem:ls-spectral-function-calculus}
设$\bm S$实对称，$f,g$定义在$\bm S$的全部特征值上。
\begin{enumerate}
 \item 对实系数多项式$p(z)=\sum_{k=0}^mc_kz^k$，按矩阵幂与加法算出的$\sum_kc_k\bm S^k$等于式~\eqref{eq:ls-symmetric-spectral-function}定义的$p(\bm S)$。
 \item $f(\bm S)+g(\bm S)=(f+g)(\bm S)$，$f(\bm S)g(\bm S)=(fg)(\bm S)$。
 \item 若$f$在全部特征值上不为零，则$f(\bm S)$可逆，且$f(\bm S)^{-1}=(1/f)(\bm S)$；特别地，$\bm S$可逆时取$f(z)=1/z$即得$\bm S^{-1}$。
\end{enumerate}
\end{lemma}
\begin{proof}
先证第二条。加法由式~\eqref{eq:ls-symmetric-spectral-function}的第二式直接得到。乘法利用谱投影的正交关系：
\[
 f(\bm S)g(\bm S)
 =\sum_\lambda\sum_\mu f(\lambda)g(\mu)\bm P_\lambda\bm P_\mu
 =\sum_\lambda f(\lambda)g(\lambda)\bm P_\lambda
 =(fg)(\bm S),
\]
中间一步用到$\lambda\neq\mu$时$\bm P_\lambda\bm P_\mu=\bm0$，以及$\bm P_\lambda^2=\bm P_\lambda$。

第一条：常数函数$1$给出$\sum_\lambda\bm P_\lambda=\bm I=\bm S^0$，恒等函数$z$给出$\sum_\lambda\lambda\bm P_\lambda=\bm S$。由乘法规则对$k$归纳，$z^k$给出$\bm S^k$；再由加法规则（数乘同理）按系数相加，得到$p(\bm S)=\sum_kc_k\bm S^k$。

第三条：由乘法规则，$f(\bm S)\,(1/f)(\bm S)=(f\cdot 1/f)(\bm S)$，即常数函数$1$作用的结果$\bm I$；两个因子的顺序交换后同理。$f(z)=1/z$时，$(1/f)(z)=z$，于是$(1/z)(\bm S)$是$\bm S$的逆。
\end{proof}

第一条也可以直接看出：连乘$\bm S^k=\bm Q\bm\Lambda\bm Q^\top\cdots\bm Q\bm\Lambda\bm Q^\top$时相邻的$\bm Q^\top\bm Q$相消，得$\bm S^k=\bm Q\bm\Lambda^k\bm Q^\top$。所以对多项式而言，沿谱作用与按矩阵运算计算是一回事；定义~\ref{def:ls-symmetric-matrix-function}只是把这个事实推广到一般的$f$。由第二条，$(\sqrt z)^2=z$给出$(\bm S^{1/2})^2=\bm S$，$\sqrt z\cdot(1/\sqrt z)=1$给出$\bm S^{1/2}\bm S^{-1/2}=\bm I$，与前面的直接验证一致。

矩阵函数沿特征方向作用，而不是对每个元素作用。下面的例子对照这两种做法。

\begin{example}[矩阵平方根与逐元素开方]
\label{ex:ls-matrix-square-root-two-by-two}
取示例~\ref{ex:ls-symmetric-spectral-two-by-two}的$\bm S=\left[\begin{smallmatrix}2&1\\1&2\end{smallmatrix}\right]=3\bm q_1\bm q_1^\top+\bm q_2\bm q_2^\top$，其中$\bm q_1=(1,1)^\top/\sqrt2$，$\bm q_2=(1,-1)^\top/\sqrt2$。按式~\eqref{eq:ls-symmetric-spectral-function}取$f(z)=\sqrt z$，
\[
 \bm S^{1/2}
 =\sqrt3\,\bm q_1\bm q_1^\top+\bm q_2\bm q_2^\top
 =\frac12\begin{bmatrix}\sqrt3+1&\sqrt3-1\\\sqrt3-1&\sqrt3+1\end{bmatrix}.
\]
计算它的平方：对角元为$\tfrac14\bigl[(\sqrt3+1)^2+(\sqrt3-1)^2\bigr]=\tfrac14\cdot8=2$，非对角元为$\tfrac14\cdot2(\sqrt3+1)(\sqrt3-1)=\tfrac14\cdot4=1$，所以$(\bm S^{1/2})^2=\bm S$。同样取$f(z)=1/\sqrt z$，得逆平方根
\[
 \bm S^{-1/2}
 =\frac1{\sqrt3}\bm q_1\bm q_1^\top+\bm q_2\bm q_2^\top
 =\frac1{2\sqrt3}\begin{bmatrix}\sqrt3+1&1-\sqrt3\\1-\sqrt3&\sqrt3+1\end{bmatrix}.
\]
逐元素开方得到的矩阵及其平方则是
\[
 \begin{bmatrix}\sqrt2&1\\1&\sqrt2\end{bmatrix}^2
 =\begin{bmatrix}3&2\sqrt2\\2\sqrt2&3\end{bmatrix}\neq\bm S .
\]
结论：沿特征方向定义的$\bm S^{1/2}$满足$(\bm S^{1/2})^2=\bm S$，逐元素开方的结果不满足所需的矩阵乘法关系。
\end{example}

固定的矩阵只有有限个特征值，所以式~\eqref{eq:ls-symmetric-spectral-function}只需要$f$在这些点上的值，不要求$f$连续，也不需要矩阵级数或微分。

还要把沿谱构造的矩阵函数与第~\ref{sec:multivariable-matrix-perturbation}节研究的一般矩阵输入映射区分开。行列式$\bm X\mapsto\det\bm X$、矩阵乘法$(\bm X,\bm Y)\mapsto\bm X\bm Y$、求逆等都以矩阵为输入，范围比$\bm S\mapsto f(\bm S)$广得多。即使研究$\bm S\mapsto f(\bm S)$本身的微分，扰动$\bm E$一般也不保持$\bm S$的特征方向，不能直接套用标量求导公式。例如$f(z)=z^2$，
\[
 (\bm S+\bm E)^2-\bm S^2=\bm S\bm E+\bm E\bm S+\bm E^2 ,
\]
一阶部分是$\bm S\bm E+\bm E\bm S$；只有$\bm S$与$\bm E$可交换时，它才化为标量公式对应的$2\bm S\bm E$。

\paragraph{实幂与对数的补充例子}

对正定$\bm S$，全部$\lambda_i>0$，任意实数$\alpha$都使$z^\alpha$在谱上有定义，于是
\[
 \bm S^\alpha=\bm Q\operatorname{diag}(\lambda_1^\alpha,\ldots,\lambda_d^\alpha)\bm Q^\top .
\]
由引理~\ref{lem:ls-spectral-function-calculus}，$\bm S^\alpha\bm S^\beta=\bm S^{\alpha+\beta}$；特别地，对正整数$m$，$(\bm S^{1/m})^m=\bm S$，即把非负伸缩分成$m$次相同的作用时，每次沿各主方向伸缩$\lambda_i^{1/m}$倍；$\bm S^{1/m}$本身对称正定，因为它的谱分解中特征值$\lambda_i^{1/m}>0$。这里只用到它是$\bm X^m=\bm S$的一个对称正定解，本节不讨论$m>2$时这类解的唯一性。若$\bm S$半正定但有零特征值，$\alpha>0$时$0^\alpha=0$，$\bm S^\alpha$仍有实值的谱定义；$\alpha<0$时$0^\alpha$没有定义，负幂不能沿用。前面平方根只要求半正定、逆平方根却要求正定，正是这一区别。

对正定$\bm S$还可以取$f(z)=\log z$，得到\term{矩阵对数}（matrix logarithm）
\[
 \log\bm S=\bm Q\operatorname{diag}(\log\lambda_1,\ldots,\log\lambda_d)\bm Q^\top .
\]
它记录的是各主方向上伸缩倍数的对数：$\log\lambda_i>0$表示沿$\bm q_i$拉长，$\log\lambda_i<0$表示压缩。上式本身就是对称矩阵$\log\bm S$的正交谱分解，特征值为$\log\lambda_i$。由定义~\ref{def:ls-symmetric-matrix-function}与基的选取无关，可以用这组$\bm Q$对$\log\bm S$作用$f(z)=e^z$（矩阵指数在本节也按式~\eqref{eq:ls-symmetric-spectral-function}沿谱定义），得$e^{\log\bm S}=\bm Q\operatorname{diag}(e^{\log\lambda_i})\bm Q^\top=\bm S$。反过来，对任意实对称$\bm T=\bm Q\operatorname{diag}(\mu_i)\bm Q^\top$，$e^{\bm T}=\bm Q\operatorname{diag}(e^{\mu_i})\bm Q^\top$是特征值全为正的对称矩阵，即正定，同理$\log e^{\bm T}=\bm T$。所以在沿谱的定义下，$\log$与$e^{(\cdot)}$在正定矩阵与实对称矩阵之间互为逆映射。$\bm S$有零特征值时$\log0$无定义，这一定义不能使用。

要区分矩阵$\log\bm S$与标量$\log\det\bm S$。由$\det\bm S=\det\bm\Lambda=\prod_i\lambda_i$及迹的循环性，
\[
 \log\det\bm S
 =\sum_{i=1}^d\log\lambda_i
 =\operatorname{tr}(\bm Q\operatorname{diag}(\log\lambda_i)\bm Q^\top)
 =\operatorname{tr}(\log\bm S).
\]
左端是一个数，$\log\bm S$是$d$阶矩阵，二者通过迹联系。对数在本节只是谱函数的又一个例子，不参与极分解的证明。第~\ref{sec:multivariable-logdet-differential}节求$\log\det$的微分时，使用的是行列式的微分与标量对数的链式法则，并不以矩阵对数的微分为前提。
